\begin{figure}[p]
\centering
\includegraphics[width=\textwidth]{supplementary/curated/image_generalization.pdf}
\caption{\textbf{The neuron-identity bottleneck generalizes across image and feedback controls.} \textbf{A}, Flattened CIFAR-10 in the shunting architecture, with matched-width scalar fallback, random rank-four, exact-path and backpropagation rules (five seeds). \textbf{B}, Raw-additive CIFAR-10 with strict scalar, neuron-specific, exact-path and backpropagation feedback (twenty paired seeds; 95\% Student-$t$ intervals). Exact path is 0.859 percentage points below neuron-specific feedback and equivalent to backpropagation within the prespecified one-point margin. \textbf{C}, Fashion-MNIST in shunting and raw-additive trees (ten paired seeds; 95\% bootstrap intervals); path resolution adds no reliable benefit. \textbf{D,E}, Dominant and sub-point contrasts under fixed random soma feedback (DFA), with fifteen paired MNIST seeds and 95\% bootstrap intervals. DFA preserves the neuron-versus-scalar benefit, but does not add a second dendritic layer or increase the dimensionality of the ten-class readout error. Percentage-point scales differ.}
\label{fig:si_image_generalization}
\label{fig:regular_regimes}
\label{fig:supp_between_within_factorial}
\end{figure}

\begin{figure}[p]
\centering
\includegraphics[width=\textwidth]{supplementary/curated/mnist_dictionary_geometry.pdf}
\caption{\textbf{Intermediate dictionaries capture more image-task credit without improving accuracy.} \textbf{A,B}, Six-rule MNIST ladder in shunting and raw-additive networks; ten fresh paired seeds per architecture. Filled/open symbols use development-selected/original common rates. Decoder-only learning freezes the initialized core. \textbf{C}, Paired accuracy differences for three profiles versus a projected common profile, and exact versus three-profile credit. \textbf{D}, Initial and trained activation-error capture by one or three profiles, averaged over 2,048 images within each seed. Projection coefficients use the exact field. Means and bars in A--D have descriptive 95\% seed-bootstrap intervals. \textbf{E--G}, Independent fifteen-seed cohort: common-state branch-gradient cosine, soma-normalized exact voltage-error magnitude by depth, and within-depth path-specific error energy. These distinguish direction, amplitude and spatial variation; exact-path cosine one is definitional. Greater path-specific variation did not yield an image accuracy benefit. Voltage- and activation-space capture are different quantities. All image cohorts have one dendritic layer and a linear readout.}
\label{fig:si_mnist_dictionary_geometry}
\label{fig:supp_image_diagnostics}
\label{fig:supp-image-ladder-controls}
\end{figure}

\begin{figure}[p]
\centering
\includegraphics[width=\textwidth]{supplementary/curated/input_coverage_depth.pdf}
\caption{\textbf{Input coverage and ordinary-task depth are distinct from credit resolution.} \textbf{A--D}, Spatial versus randomly assigned contacts in depth-four $[2,2,2,2]$ trees: example map, unique input coverage, cross-branch overlap and learning benefit under backpropagation and three local fields. Contact count is matched, but spatial maps reach more unique inputs and avoid collisions. The accuracy benefit persists under backpropagation and therefore reflects a forward sparsity prior. \textbf{E,F}, Separate raw-additive historical noise cohort with sixteen terminals: all four D4-minus-D1 accuracy contrasts are negative in 10/10 paired seeds; active contact counts are approximately fixed while parameter counts vary. Learning contrasts use means and 95\% paired-seed bootstrap intervals; map summaries use mean $\pm$ s.d. The noise-task executed generator remains unresolved and its shunting runs fail the positive-conductance input criterion. These retained additive outcomes are a depth boundary, not evidence for shunting; exclusions are enumerated in Table~\ref{tab:input_validity}.}
\label{fig:si_input_coverage_depth}
\label{fig:supp_spatial_audit}
\label{fig:supp_fixed_budget}
\end{figure}
